\paragraph{Dependence sensitivity and system-level audit.}
The dependence sensitivity keeps the original precision threshold, policy
choices and recorded outcomes. It adds fourteen-day circular calendar blocks
to the primary seven-day bootstrap, using 2,000 draws and
\texttt{PCG64 SeedSequence([20261013,cohort,14])}, where cohort is 0 for the
original panel and 1 for the external panel. Each draw concatenates sampled
blocks, wraps at the end of the observed calendar, and truncates to its
original length. The same dates are sampled across all systems, intraday
blocks, policies and information regimes. Compatible record counts supply
the denominator within each resample. Percentile endpoints use linear
interpolation at 2.5\% and 97.5\%. The fourteen-day block is long relative
to the external 31-day horizon: it is a sensitivity check, not evidence that
all temporal dependence is captured. Leave-one-system-out results retain all
records from the remaining two systems. Neither omitting one system nor
resampling days provides independent deployments or population inference.

The paired ledgers retain all four success cells, actual capacity compatibility
flags, and each system's effect. The archived fourteen-day draw indices and
the separately written checker permit direct recounting without refitting or
replanning. The checker recomputes cells with rational arithmetic and scalar
day-multiplicity sums, checks percentile interpolation independently, confirms
that the seven-day intervals are unchanged, and checks precision-gap counts
against every archived case. The full coefficient-only plans were also matched
against all 1,098 derived precision records; public replay needs only those
derived ledgers, not the original raw-dependent plans.
